%% === FORMULAS CLAVE - Practica 2: Osciloscopio y Figuras de Lissajous ===
\subsection*{Practica 2 - Osciloscopio y Lissajous}
\begin{itemize}
  \item \textbf{Lectura por rejilla:} $V_{pp}=n_V(\mathrm{V/div})$, $T=n_T(\mathrm{s/div})$ \hfill \textit{para obtener magnitudes directas en modo X--T}
  \item \textbf{Frecuencia:} $f=1/T$ \hfill \textit{a partir del periodo medido}
  \item \textbf{Voltaje de pico y eficaz:} $V_p=V_{pp}/2$, $V_{ef}=V_{pp}/(2\sqrt{2})$ \hfill \textit{para senoidal}
  \item \textbf{Incertidumbre tipo B:} $\delta x=\text{resolución}/2$ \hfill \textit{medidas directas}
  \item \textbf{Propagación general:} $\delta f=\sqrt{\sum_i\left(\frac{\partial f}{\partial x_i}\delta x_i\right)^2}$ \hfill \textit{magnitudes calculadas}
  \item \textbf{Propagación de frecuencia:} $\delta f=\delta T/T^2$ \hfill \textit{cuando $f=1/T$}
  \item \textbf{Señales en modo X--Y:} $X=A\sin(\omega_1 t)$, $Y=B\sin(\omega_2 t+\phi)$ \hfill \textit{base de Lissajous}
  \item \textbf{Circunferencia en Lissajous:} $\nu_1=\nu_2$, $A=B$, $\phi=\pm\pi/2+2k\pi$ \hfill \textit{condiciones necesarias}
\end{itemize}
